\documentclass[12pt,a4paper]{article}

\usepackage[a4paper,top=1cm,bottom=1cm,left=2cm,right=2cm]{geometry}
\usepackage{amsmath,amssymb,amsthm,mathtools}
\usepackage{xcolor}
\usepackage[most]{tcolorbox}
\usepackage{fancyhdr}
\usepackage{setspace}
\usepackage{enumitem}
\usepackage[
    unicode=true,
    bookmarks=true,
    bookmarksopen=true,
    bookmarksnumbered=true,
    bookmarksopenlevel=3,
    hidelinks
]{hyperref}
\usepackage{bookmark}

% XePersian is loaded after the other packages.
\usepackage{xepersian}

% B Mitra: local Persian fonts, with genuine bold weight.
\settextfont[
    Path=./fonts/,
    Extension=.ttf,
    UprightFont={B Mitra_0},
    BoldFont={B Mitra Bold_0}
]{B Mitra}

% Times New Roman: local Latin fonts.
\setlatintextfont[
    Path=./fonts/,
    UprightFont=times.ttf,
    BoldFont=timesbd.ttf,
    ItalicFont=timesi.ttf,
    BoldItalicFont=timesbi.ttf
]{times.ttf}

% English terms and student ID must be typeset left-to-right.
\newcommand{\en}[1]{\lr{#1}}

% B Mitra does NOT provide U+2022. Use the mathematical bullet.
% This is the only list-marker definition; do not override it later.
\setlist[itemize,1]{
    label={$\bullet$},
    rightmargin=1.2em,
    leftmargin=0pt,
    itemsep=0.5em,
    topsep=0.4em
}

% Change these values.
\newcommand{\studentname}{حامی عبادزاده سمنانی}
\newcommand{\studentid}{40552521603}
\newcommand{\coursename}{یادگیری ماشین}
\newcommand{\assignmentname}{تمرین تئوری ـ سری اول}
\newcommand{\semestername}{نیم‌سال تحصیلی ۴۰۵۱}

\definecolor{ink}{HTML}{253445}
\definecolor{accent}{HTML}{315C76}
\definecolor{softink}{HTML}{687889}
\definecolor{questionbg}{HTML}{F4F7F9}
\definecolor{rulecolor}{HTML}{DCE4EA}

\setstretch{1.24}
\setlength{\parindent}{0pt}
\setlength{\parskip}{0.48em}
\setlength{\headheight}{16pt}
\emergencystretch=2em

\pagestyle{fancy}
\fancyhf{}
\fancyfoot[C]{\small\textcolor{softink}{\thepage}}
\renewcommand{\headrulewidth}{0pt}
\renewcommand{\footrulewidth}{0pt}

\newtcolorbox{promptbox}{
    enhanced,breakable,
    colback=questionbg,
    colframe=questionbg,
    boxrule=0pt,
    arc=1.5mm,
    left=11pt,right=11pt,top=8pt,bottom=8pt,
    borderline east={2pt}{0pt}{accent},
    before skip=8pt,after skip=11pt
}

% The sectioning commands generate both numbering and clickable PDF bookmarks.
% Article class: section = question, subsection = topic,
% subsubsection = proof or a smaller conceptual part.
\setcounter{secnumdepth}{3}
\setcounter{tocdepth}{3}
\bookmarksetup{open,depth=3,numbered}

\newcommand{\answerheading}{%
    \par\vspace{0.10em}%
    {\color{accent}\bfseries پاسخ}\par\vspace{-0.2em}%
}

\begin{document}
\pdfbookmark[0]{\coursename ـ \assignmentname}{document-title}

\begin{center}
    {\small\textcolor{softink}{به نام خدا}}\par
    {\color{ink}\LARGE\bfseries \coursename}\par
    {\color{accent}\large \assignmentname}\par
    {\small\textcolor{softink}{\semestername}}\par
\end{center}

\vspace{0.85em}
{\color{accent}\hrule height 1.15pt}
\vspace{0.70em}

\noindent
\begin{minipage}[t]{0.56\textwidth}
    {\small\textcolor{softink}{نام}}\par\vspace{0.18em}
    {\bfseries \studentname}
\end{minipage}%
\begin{minipage}[t]{0.44\textwidth}
    \begin{flushleft}
        {\small\textcolor{softink}{شماره دانشجویی}}\par\vspace{0.18em}
        {\bfseries\lr{\studentid}}
    \end{flushleft}
\end{minipage}

\vspace{0.20em}
{\color{rulecolor}\hrule height 0.6pt}

% ====================== Question 1 ======================
\section{انواع یادگیری}
\begin{promptbox}
توضیح انواع یادگیری‌ها، به‌جز یادگیری تقویتی، به‌صورت مفهومی به همراه مثال.
\end{promptbox}
\answerheading

\subsection{\texorpdfstring{\en{Supervised Learning}}{Supervised Learning}}

% IMPORTANT: B Mitra supports U+06C0 but lacks U+0654.
% We therefore use the font-compatible single character «ۀ» in
% «مجموعه‌دادۀ» and «مسئلۀ» below. Avoid decomposed heh + combining hamza here;
% its combining U+0654 glyph is missing in this font.
در یادگیری نظارت‌شده، مجموعه‌دادۀ آموزشی شامل نمونه‌هایی است که
برای هر یک، برچسب یا مقدار هدف مشخصی وجود دارد.
هدف، یادگیری رابطه‌ای میان ویژگی‌های ورودی و خروجی‌های متناظر است؛
به‌گونه‌ای که مدل بتواند خروجی نمونه‌های جدید و مشاهده‌نشده را پیش‌بینی کند.

این نوع یادگیری به دو مسئلۀ اصلی تقسیم می‌شود:

\begin{itemize}
    \item \textbf{\en{Regression}:}
    در رگرسیون، متغیر هدف معمولاً پیوسته و عددی است و مدل
    تلاش می‌کند مقدار آن را برای ورودی‌های جدید تخمین بزند.

    \item \textbf{\en{Classification}:}
    در دسته‌بندی، متغیر هدف گسسته است و به مجموعه‌ای
    از کلاس‌های ازپیش‌تعریف‌شده تعلق دارد.
    هدف مدل، تعیین کلاس مناسب برای هر ورودی جدید است.
\end{itemize}

\subsubsection{تمایز مفهومی رگرسیون و دسته‌بندی}
در رگرسیون، مقدار پیش‌بینی‌شده لزوماً نباید یکی از مقادیر
مشاهده‌شده در مجموعۀ آموزشی باشد؛ برای مثال، مدل می‌تواند
مقداری میان دو مقدار آموزشی یا حتی خارج از بازۀ آن‌ها پیش‌بینی کند.
در مقابل، در دسته‌بندی، خروجی نهایی مدل باید
یکی از کلاس‌های تعریف‌شده برای مسئله باشد و نمی‌تواند
کلاس جدیدی خارج از این مجموعه ایجاد کند.

\subsubsection{مثال‌ها}
\begin{itemize}
    \item \textbf{رگرسیون:}
    پیش‌بینی قیمت یک خانه بر اساس متراژ، موقعیت مکانی
    و سال ساخت آن. قیمت پیش‌بینی‌شده می‌تواند مقداری باشد
    که دقیقاً در هیچ‌یک از نمونه‌های آموزشی وجود نداشته است.

    \item \textbf{دسته‌بندی:}
    تشخیص خوش‌خیم یا بدخیم بودن یک تومور بر اساس
    ویژگی‌های اندازه‌گیری‌شدۀ آن.
    مدل، نمونۀ جدید را در یکی از دو کلاس
    «خوش‌خیم» یا «بدخیم» قرار می‌دهد.
\end{itemize}



\subsection{\texorpdfstring{\en{Unsupervised Learning}}{Unsupervised Learning}}

در یادگیری بدون نظارت، برخلاف یادگیری نظارت‌شده، داده‌های آموزشی
فاقد برچسب یا مقدار هدف مشخص هستند. در این نوع یادگیری،
هدف اصلی کشف ساختارها، الگوها و روابط پنهان میان داده‌ها است؛
بدون آنکه خروجی مطلوب هر نمونه از پیش در اختیار مدل قرار گرفته باشد.

برخی از مسائل مهم در یادگیری بدون نظارت عبارت‌اند از:

\begin{itemize}
    \item \textbf{\en{Clustering}:}
    هدف، تقسیم داده‌ها به گروه‌هایی است که نمونه‌های درون
    هر گروه از نظر ویژگی‌هایشان به یکدیگر شباهت بیشتری
    داشته باشند. برخلاف دسته‌بندی در یادگیری نظارت‌شده،
    این گروه‌ها از پیش برچسب‌گذاری نشده‌اند و مدل
    ساختار گروه‌بندی را از خود داده‌ها کشف می‌کند.

    \item \textbf{\en{Dimensionality Reduction}}

    \item \textbf{\en{Density Estimation / Anomaly Detection}:}
    در تخمین چگالی، هدف، شناخت توزیع احتمالی داده‌ها
    و تعیین نواحی پرتراکم و کم‌تراکم فضای ویژگی‌ها است.
    در تشخیص ناهنجاری نیز تلاش می‌شود نمونه‌هایی
    شناسایی شوند که با الگوهای معمول داده‌ها
    سازگاری ندارند. کم بودن چگالی می‌تواند
    نشانه‌ای از ناهنجاری باشد، اما لزوماً
    هر نمونه در نواحی کم‌تراکم، ناهنجار نیست.
\end{itemize}

\subsubsection{مثال‌ها}

\begin{itemize}
    \item \textbf{خوشه‌بندی:}
    فرض کنید فروشگاهی اطلاعاتی از میزان خرید،
    تعداد مراجعه و نوع محصولات خریداری‌شده توسط
    مشتریان خود در اختیار دارد. بدون آنکه گروه‌بندی
    مشتریان از پیش مشخص شده باشد، می‌توان با استفاده
    از خوشه‌بندی، مشتریانی با رفتار خرید مشابه
    را در گروه‌های مشترک قرار داد.

    \item \textbf{تخمین چگالی و تشخیص ناهنجاری:}
    فرض کنید داده‌های حسگرهای یک دستگاه صنعتی
    در شرایط عملکرد عادی جمع‌آوری شده‌اند.
    با تخمین توزیع این داده‌ها می‌توان محدوده‌های
    رایج عملکرد دستگاه را شناسایی کرد.
    در نتیجه، مشاهده مقادیر بسیار نامحتمل
    یا الگوهای غیرمعمول می‌تواند نشانه‌ای
    از بروز نقص احتمالی در دستگاه باشد.
\end{itemize}



% ============== Semi-supervised Learning ==============

\subsection{\texorpdfstring{\en{Semi-supervised Learning}}{Semi-supervised Learning}}

در یادگیری نیمه‌نظارت‌شده، مجموعه‌داده آموزشی شامل
تعداد محدودی از نمونه‌های برچسب‌دار و تعداد بیشتری
از نمونه‌های بدون برچسب است.
هدف، استفاده هم‌زمان از اطلاعات نمونه‌های برچسب‌دار
و ساختار موجود در داده‌های بدون برچسب برای
بهبود عملکرد مدل است.

این روش به‌ویژه زمانی کاربرد دارد که جمع‌آوری داده
آسان باشد، اما برچسب‌گذاری آن به زمان، هزینه
یا دانش تخصصی نیاز داشته باشد.

یکی از روش‌های رایج، برچسب‌گذاری
\en{Pseudo-labeling} است؛
به این صورت که ابتدا مدل با داده‌های برچسب‌دار
آموزش می‌بیند و سپس برای نمونه‌های بدون برچسب
پیش‌بینی انجام می‌دهد. پیش‌بینی‌هایی که مدل
از صحت آن‌ها اطمینان بالایی دارد، به‌عنوان
برچسب‌های موقت در ادامه آموزش استفاده می‌شوند.

البته همه روش‌های نیمه‌نظارت‌شده مبتنی بر
تولید برچسب نیستند؛ برای مثال، برخی روش‌ها
از شباهت میان نمونه‌ها یا سازگاری پیش‌بینی‌های
مدل در برابر تغییرات کوچک ورودی استفاده می‌کنند.

\subsubsection{تمایز با یادگیری نظارت‌شده}

در یادگیری نظارت‌شده، تابع یادگیری
با استفاده از نمونه‌های برچسب‌دار آموزش می‌بیند؛
اما در یادگیری نیمه‌نظارت‌شده، داده‌های بدون برچسب
نیز در فرایند آموزش نقش دارند.
بنابراین تفاوت اصلی در نوع خروجی مدل نیست،
بلکه در اطلاعاتی است که هنگام آموزش
مورد استفاده قرار می‌گیرد.

همچنین برچسب‌های تولیدشده توسط خود مدل
لزوماً صحیح نیستند و نباید با برچسب‌های
واقعی یکسان در نظر گرفته شوند.

\subsubsection{مثال}

فرض کنید بیمارستانی یک میلیون تصویر رادیولوژی
از قفسه سینه در اختیار دارد، اما تنها دو هزار تصویر
توسط متخصصان برچسب‌گذاری شده‌اند.
می‌توان ابتدا یک مدل دسته‌بندی را با تصاویر
برچسب‌دار آموزش داد و سپس با بهره‌گیری
از اطلاعات تصاویر بدون برچسب، عملکرد آن
را بهبود بخشید؛ بدون آنکه لازم باشد
تمام تصاویر به‌صورت دستی برچسب‌گذاری شوند.


% ============== Self-supervised Learning ==============

\subsection{\texorpdfstring{\en{Self-supervised Learning}}{Self-supervised Learning}}

در یادگیری خودنظارت‌شده، برخلاف یادگیری
نظارت‌شده، نیازی به برچسب‌گذاری دستی
داده‌ها برای وظیفه آموزشی اولیه نیست.
در عوض، با استفاده از ساختار خود داده‌ها،
یک مسئله پیش‌بینی تعریف می‌شود که
اطلاعات لازم برای تعیین پاسخ صحیح آن
از خود داده قابل استخراج است.

به عبارت دیگر، مدل مسئله‌ای را حل می‌کند
که پاسخ آن از پیش در داده‌ها وجود دارد،
اما بخشی از اطلاعات ورودی از مدل پنهان
شده یا به شکل دیگری در اختیار آن قرار می‌گیرد.

هدف معمولاً یادگیری بازنمایی‌های مفیدی
از داده‌ها است که بتوان از آن‌ها
در مسائل دیگر نیز استفاده کرد.

\subsubsection{مثال}

در آموزش مدل‌های زبانی، می‌توان بخشی
از یک جمله را در اختیار مدل قرار داد
و از آن خواست کلمه بعدی را پیش‌بینی کند.

در این حالت، کلمه صحیح از خود متن آموزشی
استخراج می‌شود و نیازی نیست انسانی
برای هر جمله برچسب جداگانه‌ای تعیین کند.

مدل با حل مکرر این مسئله، الگوهای
زبانی و روابط میان کلمات را یاد می‌گیرد.
از دانش حاصل می‌توان در وظایف دیگری
مانند دسته‌بندی متن یا پاسخ‌گویی
به پرسش‌ها استفاده کرد.

% ====================== Question 2 ======================

\section{تعداد پارامترهای چندجمله‌ای}
\begin{promptbox}
تعداد پارامتر‌های چندجمله‌ای درجه $M$ در فضای $D$ بعدی و تحلیل رفتار آن با افزایش بعد.
\end{promptbox}
\answerheading

% ==========================================
% پاسخ پرسش دوم: تعداد پارامترهای چندجمله‌ای
% ==========================================

\subsection{فرمول‌بندی چندجمله‌ای در فضای چندبعدی}

فرض کنید $D\geq1$ و $M\geq0$ اعداد صحیح باشند و
بردار ورودی
$\mathbf{x}=(x_1,\ldots,x_D)\in\mathbb{R}^D$
باشد. یک چندجمله‌ای عمومی با \textbf{درجه کل حداکثر $M$}
را می‌توان به صورت

\[
f(\mathbf{x})=
\sum_{\substack{\boldsymbol{\alpha}\in\mathbb{Z}_{\geq0}^{D}\\
                 |\boldsymbol{\alpha}|\leq M}}
w_{\boldsymbol{\alpha}}\,\mathbf{x}^{\boldsymbol{\alpha}}
\]

نوشت، که در آن

\[
\boldsymbol{\alpha}=(\alpha_1,\ldots,\alpha_D),\qquad
|\boldsymbol{\alpha}|=\sum_{j=1}^{D}\alpha_j,\qquad
\mathbf{x}^{\boldsymbol{\alpha}}=\prod_{j=1}^{D}x_j^{\alpha_j}.
\]

تمام توان‌ها اعداد صحیح نامنفی هستند؛ زیرا در تعریف
چندجمله‌ای، توان منفی وجود ندارد. هر بردار توان
$\boldsymbol{\alpha}$ یک \textbf{تک‌جمله‌ای متمایز}
را مشخص می‌کند. برای مثال، $x_1x_2$ و $x_2x_1$
هر دو به بردار $\boldsymbol{\alpha}=(1,1)$ مربوط‌اند
و نباید دو بار شمرده شوند.

\subsection{قضیه: تک‌جمله‌ای‌های متمایز یک پایه تشکیل می‌دهند}

فضای برداری چندجمله‌ای‌های حقیقی با درجه کل حداکثر $M$
را چنین تعریف می‌کنیم:

\[
\mathcal{P}_{D,M}
=\bigl\{p\in\mathbb{R}[x_1,\ldots,x_D]:\deg p\leq M\bigr\}.
\]

\textbf{قضیه.}
مجموعه

\[
\mathcal{B}_{D,M}
=\left\{
\mathbf{x}^{\boldsymbol{\alpha}}:
\boldsymbol{\alpha}\in\mathbb{Z}_{\geq0}^{D},\;
|\boldsymbol{\alpha}|\leq M
\right\}
\]

یک پایه برای $\mathcal{P}_{D,M}$ است.

\subsubsection{اثبات پوشا بودن}

طبق تعریف چندجمله‌ای، هر عضو $\mathcal{P}_{D,M}$
ترکیبی خطی و متناهی از تک‌جمله‌ای‌های مجموعه
$\mathcal{B}_{D,M}$ است؛ پس

\[
\operatorname{span}(\mathcal{B}_{D,M})
=\mathcal{P}_{D,M}.
\]

\subsubsection{اثبات استقلال خطی با استقرا}

نشان می‌دهیم اگر یک ترکیب خطی از تک‌جمله‌ای‌های
متمایز روی تمام $\mathbb{R}^{D}$ صفر باشد،
تمام ضرایب آن صفرند. اثبات با استقرا بر تعداد
متغیرها $D$ انجام می‌شود.

\textbf{پایه استقرا:} $(D=1)$.
اگر چندجمله‌ای تک‌متغیره

\[
p(x)=\sum_{k=0}^{M}c_kx^k
\]

برای تمام $x\in\mathbb{R}$ صفر باشد،
بی‌نهایت ریشه دارد. از آنجا که یک چندجمله‌ای
ناصفر از درجه حداکثر $M$ بیش از $M$ ریشه متمایز
ندارد، باید همه ضرایب $c_k$ صفر باشند.

\textbf{گام استقرا.}
فرض کنید ادعا برای چندجمله‌ای‌های $D-1$ متغیره
درست باشد. اگر چندجمله‌ای $D$ متغیره
$p(\mathbf{x})$ برای تمام $\mathbf{x}\in\mathbb{R}^{D}$
صفر باشد، آن را برحسب توان‌های $x_D$ می‌نویسیم:

\[
p(x_1,\ldots,x_D)
=\sum_{k=0}^{M}x_D^k\,
q_k(x_1,\ldots,x_{D-1}),
\]

که هر $q_k$ چندجمله‌ای با حداکثر $D-1$ متغیر
و درجه کل حداکثر $M-k$ است.
هر مقدار ثابت دلخواه برای
$(x_1,\ldots,x_{D-1})$ در نظر بگیرید.
عبارت بالا نسبت به $x_D$ یک چندجمله‌ای تک‌متغیره
است که برای \textbf{تمام} $x_D\in\mathbb{R}$ صفر می‌شود.
پس طبق حالت پایه، تمام ضرایب آن باید صفر باشند:

\[
q_k(x_1,\ldots,x_{D-1})=0,
\qquad k=0,\ldots,M.
\]

چون مقادیر $(x_1,\ldots,x_{D-1})$ دلخواه بودند،
هر $q_k$ روی کل $\mathbb{R}^{D-1}$ همانی صفر است.
با فرض استقرا، همه ضرایب تک‌جمله‌ای‌های هر $q_k$
صفرند؛ بنابراین تمام ضرایب $p$ نیز صفر هستند.

پس تک‌جمله‌ای‌های متمایز مستقل خطی‌اند و به همراه
خاصیت پوشا بودن، نتیجه می‌گیریم

\begin{center}
\fbox{\strut
    $\mathcal{B}_{D,M}$
    یک پایه برای
    $\mathcal{P}_{D,M}$
    است.
}
\end{center}

\subsection{یکتایی نمایش و نبود پارامترهای تکراری}

اکنون باید نشان دهیم که شمارش عناصر این پایه
دقیقاً تعداد \textbf{پارامترهای مستقل} را می‌دهد.
فرض کنید یک چندجمله‌ای، دو نمایش با ضرایب متفاوت
روی پایه $\mathcal{B}_{D,M}$ داشته باشد:

\[
\sum_{\boldsymbol{\alpha}}a_{\boldsymbol{\alpha}}
\mathbf{x}^{\boldsymbol{\alpha}}
=
\sum_{\boldsymbol{\alpha}}b_{\boldsymbol{\alpha}}
\mathbf{x}^{\boldsymbol{\alpha}},
\qquad\forall\mathbf{x}\in\mathbb{R}^{D},
\]

که هر دو مجموع روی بردارهای توان مجاز گرفته شده‌اند.
با تفریق دو طرف خواهیم داشت:

\[
\sum_{\boldsymbol{\alpha}}
(a_{\boldsymbol{\alpha}}-b_{\boldsymbol{\alpha}})
\mathbf{x}^{\boldsymbol{\alpha}}=0.
\]

از استقلال خطی اعضای پایه نتیجه می‌شود

\begin{center}
\fbox{\strut
    $a_{\boldsymbol{\alpha}}=b_{\boldsymbol{\alpha}}$
    \quad برای هر
    $\boldsymbol{\alpha}$
}
\end{center}

بنابراین هر چندجمله‌ای دقیقاً یک بردار ضرایب
در این پایه دارد و هیچ ضریب زائدی شمرده نمی‌شود.
از سوی دیگر، چون پایه کل فضا را پوشش می‌دهد،
هیچ چندجمله‌ای ممکنی نیز از شمارش حذف نشده است.
در نتیجه

\[
\boxed{
P(D,M)=\dim\mathcal{P}_{D,M}
=|\mathcal{B}_{D,M}|
=\#\!\left\{\boldsymbol{\alpha}\in\mathbb{Z}_{\geq0}^{D}:
|\boldsymbol{\alpha}|\leq M\right\}.
}
\]

اکنون کافی است تعداد بردارهای توان مجاز را بشماریم.


\subsection{شمارش تعداد کل پارامترهای مستقل}

بر اساس قضیۀ پایه بودن تک‌جمله‌ای‌های متمایز
و یکتایی نمایش چندجمله‌ای، تعداد پارامترهای مستقل
برابر با تعداد بردارهای توان
$\boldsymbol{\alpha}\in\mathbb{Z}_{\geq 0}^{D}$
است که شرط زیر را ارضا می‌کنند:

\[
\alpha_1+\alpha_2+\cdots+\alpha_D\leq M.
\]

برای شمارش این بردارها، یک متغیر کمکی
$\alpha_0$ تعریف می‌کنیم:

\[
\alpha_0=M-\sum_{i=1}^{D}\alpha_i.
\]

از آنجا که مجموع توان‌ها حداکثر $M$ است،
داریم $\alpha_0\geq 0$.
بنابراین شرط اولیه معادل معادلۀ زیر خواهد بود:

\[
\alpha_0+\alpha_1+\cdots+\alpha_D=M,
\qquad
\alpha_i\in\mathbb{Z}_{\geq0}
\quad (i=0,\ldots,D).
\]

این تبدیل یک تناظر یک‌به‌یک میان
بردارهای توان مجاز و جواب‌های معادلۀ جدید
برقرار می‌کند؛ زیرا به ازای هر
$(\alpha_1,\ldots,\alpha_D)$،
مقدار $\alpha_0$ به‌طور یکتا تعیین می‌شود.
از طرف دیگر، حذف $\alpha_0$ از هر جواب
معادلۀ جدید، یک بردار توان مجاز به دست می‌دهد.

در نتیجه، کافی است تعداد جواب‌های صحیح
نامنفی معادلۀ اخیر را محاسبه کنیم.

\par\medskip
\noindent
{\bfseries استفاده از اصل ستاره و میله
\en{(Stars and Bars)}}
\par\smallskip

طبق این اصل، تعداد جواب‌های صحیح نامنفی معادلۀ

\[
z_1+z_2+\cdots+z_k=n
\]

برابر است با

\[
\binom{n+k-1}{k-1}.
\]

برای اثبات این رابطه، $n$ ستاره را
در یک ردیف قرار می‌دهیم و با استفاده از
$k-1$ میله، آن‌ها را به $k$ گروه تقسیم می‌کنیم.
تعداد ستاره‌های گروه $i$ مقدار $z_i$ را
مشخص می‌کند. گروه‌های خالی نیز مجاز هستند؛
بنابراین میله‌ها می‌توانند در کنار یکدیگر
یا در ابتدا و انتهای چیدمان قرار گیرند.

هر جواب معادله دقیقاً با یک چیدمان
از $n$ ستاره و $k-1$ میله متناظر است.
در مجموع $n+k-1$ جایگاه داریم که
انتخاب $k-1$ جایگاه برای میله‌ها،
تعداد چیدمان‌های ممکن را تعیین می‌کند.

اکنون در مسئلۀ ما، تعداد متغیرها
$k=D+1$ و مجموع آن‌ها $n=M$ است.
پس تعداد جواب‌ها برابر خواهد بود با

\[
\begin{aligned}
P(D,M)
&=\binom{M+(D+1)-1}{(D+1)-1}\\[4pt]
&=\binom{D+M}{D}\\[4pt]
&=\binom{D+M}{M}.
\end{aligned}
\]

بنابراین تعداد پارامترهای مستقل
یک چندجمله‌ای عمومی از درجۀ حداکثر
$M$ در فضای $D$بعدی برابر است با

\[
\boxed{
P(D,M)=\binom{D+M}{M}
=\frac{(D+M)!}{D!\,M!}.
}
\]

لازم به ذکر است که تک‌جمله‌ای ثابت
نیز در این شمارش لحاظ شده است؛
زیرا بردار توان
$\boldsymbol{\alpha}=\mathbf{0}$
شرط درجه را ارضا می‌کند و متناظر با
ضریب ثابت چندجمله‌ای است.

% ==========================================
% تحلیل رشد تعداد پارامترها
% ==========================================

\subsection{تحلیل رشد تعداد پارامترها با افزایش بعد}

برای بررسی ارتباط تعداد پارامترها با بعد
فضای ورودی، باید میان حالت‌های مختلف
رشد $D$ و $M$ تمایز قائل شویم.


\subsubsection{حالت اول: درجه \texorpdfstring{$M$}{M} ثابت و بعد \texorpdfstring{$D$}{D} متغیر}

رابطۀ به‌دست‌آمده را می‌توان به شکل زیر بازنویسی کرد:

\[
\begin{aligned}
P(D,M)
&=\frac{(D+1)(D+2)\cdots(D+M)}{M!}\\[4pt]
&=\frac{D^M}{M!}
\prod_{j=1}^{M}\left(1+\frac{j}{D}\right).
\end{aligned}
\]

از آنجا که $M$ ثابت است، تعداد عوامل ضرب
متناهی و مستقل از $D$ است.
بنابراین با استفاده از قضیۀ حد حاصل‌ضرب داریم:

\[
\begin{aligned}
\lim_{D\to\infty}
\frac{P(D,M)}{D^M/M!}
&=
\lim_{D\to\infty}
\prod_{j=1}^{M}\left(1+\frac{j}{D}\right)\\[6pt]
&=
\prod_{j=1}^{M}
\lim_{D\to\infty}\left(1+\frac{j}{D}\right)\\[6pt]
&=
\prod_{j=1}^{M}1=1.
\end{aligned}
\]

طبق تعریف هم‌ارزی مجانبی، نتیجه می‌گیریم:

\[
\boxed{
P(D,M)\sim\frac{D^M}{M!}
\qquad (D\to\infty,\ M\text{ ثابت}).
}
\]

پس برای درجۀ ثابت $M$، تعداد پارامترها
نسبت به بعد $D$ رشد چندجمله‌ای از مرتبه
$D^M$ دارد.


\subsubsection{حالت دوم: افزایش هم‌زمان درجه و بعد}

اگر درجه چندجمله‌ای نیز متناسب با بعد
افزایش یابد، ممکن است رشد پارامترها
نمایی شود.

برای اثبات این موضوع، حالت خاص
$M=D$ را در نظر می‌گیریم:

\[
P(D,D)=\binom{2D}{D}
=\frac{(2D)!}{(D!)^2}.
\]

با استفاده از تقریب استرلینگ

\[
n!\sim\sqrt{2\pi n}
\left(\frac{n}{e}\right)^n,
\]

خواهیم داشت:

\[
\begin{aligned}
P(D,D)
&\sim
\frac{
\sqrt{4\pi D}(2D/e)^{2D}
}{
2\pi D(D/e)^{2D}
}\\[4pt]
&=
\frac{4^D}{\sqrt{\pi D}}.
\end{aligned}
\]

پس

\[
\boxed{
P(D,D)\sim\frac{4^D}{\sqrt{\pi D}}.
}
\]

از آنجا که عامل اصلی رشد، $4^D$ است،
در این حالت تعداد پارامترها با افزایش
بعد، رشد نمایی دارد.

بنابراین رشد نمایی برای چندجمله‌ای‌های
با درجه کلی حداکثر $M$، یک نتیجه
بدون قید نیست؛ بلکه به نحوه تغییر
$M$ نسبت به $D$ وابسته است.

% ==========================================
% ارتباط با تعداد نمونه‌های آموزشی
% ==========================================

\subsection{ارتباط تعداد پارامترها با تعداد نمونه‌های آموزشی}

اکنون فرض کنید مجموعه آموزشی شامل
$N$ نمونه باشد:

\[
\mathcal{D}
=\{(\mathbf{x}_i,y_i)\}_{i=1}^{N}.
\]

چندجمله‌ای مورد نظر، اگرچه نسبت به
متغیرهای ورودی غیرخطی است، نسبت به
ضرایب خود یک مدل خطی محسوب می‌شود.
بنابراین می‌توان آن را به صورت

\[
f(\mathbf{x})=
\boldsymbol{\theta}^{T}\boldsymbol{\phi}(\mathbf{x})
\]

نوشت، که در آن
$\boldsymbol{\phi}(\mathbf{x})$
بردار تمام تک‌جمله‌ای‌های درجه حداکثر
$M$ و
$\boldsymbol{\theta}\in\mathbb{R}^{P}$
بردار ضرایب مدل است؛ به‌طوری‌که

\[
P=\binom{D+M}{M}.
\]

با قرار دادن تمام نمونه‌های آموزشی
در کنار یکدیگر، در حالت بدون نویز
دستگاه معادلات زیر حاصل می‌شود:

\[
\mathbf{y}=\boldsymbol{\Phi}\boldsymbol{\theta},
\]

که در آن ماتریس طراحی
$\boldsymbol{\Phi}$
دارای ابعاد

\[
\boldsymbol{\Phi}\in\mathbb{R}^{N\times P}
\]

است.

برای آنکه بتوان ضرایب مدل را به‌صورت
یکتا از این دستگاه تعیین کرد،
لازم و کافی است ماتریس طراحی
رتبه ستونی کامل داشته باشد؛ یعنی

\[
\operatorname{rank}(\boldsymbol{\Phi})=P.
\]

اما طبق ویژگی رتبه ماتریس‌ها،

\[
\operatorname{rank}(\boldsymbol{\Phi})
\leq\min(N,P).
\]

بنابراین شرط لازم برای شناسایی یکتای
تمام ضرایب عبارت است از

\[
\boxed{
N\geq P=\binom{D+M}{M}.
}
\]

اگر $N<P$ باشد، طبق قضیه رتبه و پوچی،
فضای پوچ ماتریس طراحی دارای بعد مثبت است:

\[
\dim\ker(\boldsymbol{\Phi})
=P-\operatorname{rank}(\boldsymbol{\Phi})
\geq P-N>0.
\]

در نتیجه، بردار ناصفر
$\mathbf{v}$
وجود دارد که

\[
\boldsymbol{\Phi}\mathbf{v}=\mathbf{0}.
\]

پس اگر $\boldsymbol{\theta}$ یک جواب
دستگاه باشد، برای هر عدد حقیقی $t$،
بردار

\[
\boldsymbol{\theta}'=
\boldsymbol{\theta}+t\mathbf{v}
\]

نیز همان خروجی‌های آموزشی را تولید می‌کند؛ زیرا

\[
\boldsymbol{\Phi}\boldsymbol{\theta}'
=
\boldsymbol{\Phi}\boldsymbol{\theta}
+t\boldsymbol{\Phi}\mathbf{v}
=\mathbf{y}.
\]

بنابراین، بدون فرض‌های اضافی،
وقتی تعداد نمونه‌ها از تعداد پارامترها
کمتر باشد، داده‌های آموزشی برای
تعیین یکتای ضرایب کافی نیستند.

البته شرط $N\geq P$ به‌تنهایی کافی نیست؛
زیرا نمونه‌ها باید به‌گونه‌ای باشند
که ماتریس طراحی واقعاً رتبه کامل داشته باشد.

همچنین حتی شناسایی یکتای ضرایب در
حالت بدون نویز، تضمینی برای تعمیم
مناسب مدل در حضور نویز یا روی
داده‌های مشاهده‌نشده نیست.
تعمیم به عواملی مانند توزیع داده‌ها،
نویز، پیچیدگی مدل و منظم‌سازی نیز وابسته است.

% ==========================================
% ارتباط با نفرین ابعاد
% ==========================================

\subsection{نتیجه‌گیری و ارتباط با نفرین ابعاد}

نتیجه اصلی این است که با افزایش بعد
فضای ورودی، تعداد تک‌جمله‌ای‌های ممکن
و در نتیجه تعداد پارامترهای مدل افزایش
می‌یابد. بنابراین، برای شناسایی
ضرایب یک مدل چندجمله‌ای عمومی،
به تعداد بیشتری از نمونه‌های آموزشی
مستقل و اطلاعاتی نیاز خواهیم داشت.

برای درجه ثابت $M$، این افزایش از
مرتبه $D^M$ است؛ اما اگر درجه نیز
متناسب با بعد افزایش پیدا کند،
رشد تعداد پارامترها می‌تواند نمایی باشد.

علاوه بر این، نفرین ابعاد از دیدگاه
پوشش فضای ورودی نیز قابل مشاهده است.
فرض کنید هر یک از $D$ بعد را به $k$
بازه تقسیم کنیم. در این صورت،
تعداد نواحی حاصل برابر است با

\[
\boxed{k^D}.
\]

اگر بخواهیم به‌طور متوسط در هر ناحیه
$r$ نمونه داشته باشیم، تعداد کل
نمونه‌های مورد نیاز تقریباً برابر خواهد بود با

\[
N\approx r\,k^D.
\]

این رابطه نشان می‌دهد که حفظ تراکم
یکسان نمونه‌ها در تقسیم‌بندی فضای ورودی،
با افزایش بعد، به تعداد نمونه‌هایی
با رشد نمایی نیاز دارد.

البته این استدلال هندسی با شمارش
پارامترهای چندجمله‌ای متفاوت است
و یک شرط عمومی برای تعداد نمونه‌های
مورد نیاز تمام الگوریتم‌های یادگیری نیست.

در مجموع، افزایش بعد می‌تواند هم
پیچیدگی نمایش و محاسبات مدل را افزایش دهد
و هم یادگیری قابل اعتماد از داده‌های محدود
را دشوارتر کند. این پدیده یکی از
جلوه‌های مهم \textbf{نفرین ابعاد}
\en{(Curse of Dimensionality)}
در یادگیری ماشین است.


% ====================== Question 3 ======================

% ==========================================
% پرسش سوم: VC Dimension
% ==========================================

\section{بعد وی‌سی و ظرفیت تعمیم}

\begin{promptbox}
بعد وی‌سی \en{(VC Dimension)} چیست؟
این مفهوم را با جزئیات توضیح دهید.
\end{promptbox}

\answerheading

% ==========================================
% 1. Introduction and formal definitions
% ==========================================

\subsection{تعریف ریاضی بعد وی‌سی}

در یادگیری ماشین، یکی از مسائل اساسی،
سنجش ظرفیت یا توان بیان یک خانواده از مدل‌ها است.
برای مثال، یک مدل ممکن است بتواند داده‌های آموزشی
را کاملاً یاد بگیرد، اما روی داده‌های جدید عملکرد
مناسبی نداشته باشد.

بعد وی‌سی \en{(Vapnik--Chervonenkis Dimension)}
معیاری ترکیبیاتی برای سنجش ظرفیت یک
\textbf{کلاس فرضیه}
\en{(Hypothesis Class)}
در مسائل دسته‌بندی دودویی است.
این معیار مشخص می‌کند یک کلاس از مدل‌ها تا چه
اندازه توانایی بازتولید برچسب‌گذاری‌های دلخواه
روی مجموعه‌های متناهی از نقاط را دارد.

فرض کنید $\mathcal{X}$ فضای ورودی و

\[
\mathcal{H}
\subseteq
\{h:\mathcal{X}\rightarrow\{0,1\}\}
\]

یک کلاس فرضیه باشد.

همچنین مجموعه‌ای از $n$ نقطه متمایز را
به صورت زیر در نظر بگیرید:

\[
S=\{\mathbf{x}_1,\mathbf{x}_2,\ldots,\mathbf{x}_n\}
\subseteq\mathcal{X}.
\]

هر فرضیه $h\in\mathcal{H}$،
یک بردار برچسب روی این نقاط تولید می‌کند:

\[
\bigl(
h(\mathbf{x}_1),\ldots,h(\mathbf{x}_n)
\bigr)
\in\{0,1\}^{n}.
\]

مجموعه تمام برچسب‌گذاری‌های قابل تولید توسط
کلاس فرضیه روی $S$ را چنین تعریف می‌کنیم:

\[
\mathcal{H}(S)=
\left\{
\bigl(h(\mathbf{x}_1),\ldots,h(\mathbf{x}_n)\bigr)
\;:\;
h\in\mathcal{H}
\right\}.
\]

از آنجا که هر نقطه می‌تواند یکی از دو برچسب
صفر یا یک را دریافت کند، تعداد تمام
برچسب‌گذاری‌های ممکن برابر $2^n$ است.
بنابراین همواره داریم:

\[
|\mathcal{H}(S)|\leq 2^n.
\]

\subsubsection{شکستن کامل یک مجموعه \texorpdfstring{\en{Shattering}}{Shattering}}

می‌گوییم کلاس فرضیه $\mathcal{H}$
مجموعه $S$ را
\textbf{می‌شکند}
اگر بتواند تمام $2^n$
برچسب‌گذاری ممکن روی نقاط آن را محقق کند؛ یعنی

\[
\boxed{
\mathcal{H}(S)=\{0,1\}^{n}.
}
\]

به بیان دقیق‌تر،

\[
\boxed{
\forall \mathbf{y}\in\{0,1\}^{n},
\quad
\exists h\in\mathcal{H}:
\quad
h(\mathbf{x}_i)=y_i
\quad
\forall i\in\{1,\ldots,n\}.
}
\]

نکته مهم در این تعریف، ترتیب سورها است:
برای هر برچسب‌گذاری دلخواه، باید
\textit{حداقل یک فرضیه} وجود داشته باشد
که آن برچسب‌گذاری را محقق کند.
این فرضیه می‌تواند برای برچسب‌گذاری‌های
مختلف، متفاوت باشد.

\subsubsection{تعریف بعد وی‌سی}

بعد وی‌سی کلاس فرضیه $\mathcal{H}$،
بزرگ‌ترین تعداد نقاطی است که
\textit{حداقل یک مجموعه} با آن تعداد نقطه
توسط $\mathcal{H}$ شکسته شود.

تعریف رسمی آن عبارت است از:

\[
\boxed{
\operatorname{VCdim}(\mathcal{H})
=
\sup
\left\{
|S|:
S\subseteq\mathcal{X},
\;
S \text{ توسط }\mathcal{H}
\text{ شکسته شود}
\right\}.
}
\]

اگر مجموعه‌هایی با اندازه دلخواه بزرگ
قابل شکستن باشند، بعد وی‌سی نامتناهی است.

از این تعریف، دو اصل مهم برای اثبات
بعد وی‌سی نتیجه می‌شود:

\begin{itemize}

\item برای اثبات
$\operatorname{VCdim}(\mathcal{H})\geq d$،
کافی است یک مجموعه $d$ نقطه‌ای پیدا کنیم
که توسط $\mathcal{H}$ شکسته شود.

\item برای اثبات
$\operatorname{VCdim}(\mathcal{H})\leq d$،
باید نشان دهیم هیچ مجموعه‌ای از
$d+1$ نقطه متمایز قابل شکستن نیست.

\end{itemize}

بنابراین وجود یک مجموعه غیرقابل‌شکستن
از $d+1$ نقطه، به‌تنهایی برای اثبات
کران بالا کافی نیست؛ بلکه باید
تمام مجموعه‌های ممکن با این اندازه
را در نظر گرفت.


% ==========================================
% 2. Example: threshold classifiers
% ==========================================

\subsection{مثال اول: دسته‌بندهای آستانه‌ای}

کلاس فرضیه زیر را روی محور حقیقی
در نظر بگیرید:

\[
\mathcal{H}_{\mathrm{thr}}
=
\left\{
h_t(x)=\mathbb{I}\{x\geq t\}
:
t\in\mathbb{R}
\right\},
\]

که در آن $\mathbb{I}$ تابع شاخص است.

ادعا می‌کنیم:

\[
\boxed{
\operatorname{VCdim}(\mathcal{H}_{\mathrm{thr}})=1.
}
\]

\subsubsection{اثبات کران پایین}
یک نقطه دلخواه $x_1$ را در نظر بگیرید.

با انتخاب $t>x_1$،
برچسب صفر و با انتخاب $t\leq x_1$،
برچسب یک تولید می‌شود.

پس هر دو برچسب‌گذاری ممکن برای این
نقطه قابل تحقق است و در نتیجه

\[
\operatorname{VCdim}(\mathcal{H}_{\mathrm{thr}})
\geq 1.
\]

\subsubsection{اثبات کران بالا}
دو نقطه متمایز دلخواه $x_1<x_2$
را در نظر بگیرید.

فرض کنید بخواهیم برچسب‌های زیر را
به آن‌ها اختصاص دهیم:

\[
h_t(x_1)=1,
\qquad
h_t(x_2)=0.
\]

شرط اول مستلزم $x_1\geq t$ است؛
اما چون $x_2>x_1$،
خواهیم داشت $x_2\geq t$،
که با شرط دوم تناقض دارد.

بنابراین هیچ مجموعه دو نقطه‌ای
توسط این کلاس شکسته نمی‌شود و

\[
\operatorname{VCdim}(\mathcal{H}_{\mathrm{thr}})
\leq 1.
\]

از ترکیب دو کران نتیجه می‌گیریم:

\[
\boxed{
\operatorname{VCdim}(\mathcal{H}_{\mathrm{thr}})=1.
}
\]

این مثال نشان می‌دهد که حتی اگر
تعداد فرضیه‌های یک کلاس نامتناهی باشد،
بعد وی‌سی آن می‌تواند محدود باشد؛
زیرا در اینجا پارامتر $t$ هر مقدار
حقیقی را می‌پذیرد.


% ==========================================
% 3. Affine linear classifiers
% ==========================================

\subsection{مثال دوم: دسته‌بندهای خطی در فضای چندبعدی}

کلاس تمام دسته‌بندهای خطی آفین
در فضای $\mathbb{R}^{D}$ را در نظر بگیرید:

\[
\mathcal{H}_{\mathrm{lin}}
=
\left\{
h_{\mathbf{w},b}(\mathbf{x})
=
\mathbb{I}
\left\{
\mathbf{w}^{T}\mathbf{x}+b\geq0
\right\}
:
\mathbf{w}\in\mathbb{R}^{D},
\;
b\in\mathbb{R}
\right\}.
\]

برای این کلاس داریم:

\[
\boxed{
\operatorname{VCdim}(\mathcal{H}_{\mathrm{lin}})
=D+1.
}
\]

\subsubsection{اثبات کران پایین}
نقاط زیر را انتخاب می‌کنیم:

\[
\mathbf{x}_0=\mathbf{0},
\qquad
\mathbf{x}_i=\mathbf{e}_i,
\quad i=1,\ldots,D,
\]

که در آن $\mathbf{e}_i$ بردار پایه
استاندارد در راستای مختصه $i$ است.

برای هر برچسب‌گذاری دلخواه
$s_0,\ldots,s_D\in\{-1,+1\}$،
پارامترهای مدل را به شکل زیر انتخاب می‌کنیم:

\[
b=s_0,
\qquad
w_i=s_i-s_0.
\]

در این صورت،

\[
\mathbf{w}^{T}\mathbf{x}_0+b=s_0
\]

و برای هر $i=1,\ldots,D$ داریم:

\[
\mathbf{w}^{T}\mathbf{e}_i+b
=w_i+b=s_i.
\]

پس علامت تابع تصمیم در تمام نقاط،
دقیقاً با برچسب‌های دلخواه منطبق است.

بنابراین $D+1$ نقطه قابل شکستن هستند و

\[
\operatorname{VCdim}(\mathcal{H}_{\mathrm{lin}})
\geq D+1.
\]

\subsubsection{اثبات کران بالا}
اکنون هر مجموعه دلخواه از
$D+2$ نقطه را در نظر بگیرید.

بردارهای افزوده زیر متعلق به
$\mathbb{R}^{D+1}$ هستند:

\[
\widetilde{\mathbf{x}}_i
=
\begin{pmatrix}
\mathbf{x}_i\\
1
\end{pmatrix},
\qquad i=1,\ldots,D+2.
\]

چون تعداد این بردارها از بعد فضای
آن‌ها بیشتر است، وابسته خطی هستند.
بنابراین ضرایبی
$a_1,\ldots,a_{D+2}$،
نه همگی صفر، وجود دارند که

\[
\sum_{i=1}^{D+2}
a_i\widetilde{\mathbf{x}}_i
=\mathbf{0}.
\]

در نتیجه،

\[
\sum_i a_i\mathbf{x}_i=\mathbf{0},
\qquad
\sum_i a_i=0.
\]

از آنجا که ضرایب همگی صفر نیستند
و مجموع آن‌ها صفر است، دست‌کم
یک ضریب مثبت و یک ضریب منفی وجود دارد.

به نقاط دارای $a_i>0$ برچسب یک،
و به نقاط دارای $a_i<0$ برچسب صفر
اختصاص می‌دهیم.
برچسب نقاطی که $a_i=0$ دارند دلخواه است.

اگر این برچسب‌گذاری با تابع آفین
$f(\mathbf{x})=\mathbf{w}^{T}\mathbf{x}+b$
قابل تحقق باشد، آنگاه برای
$a_i>0$ خواهیم داشت
$a_i f(\mathbf{x}_i)\geq0$
و برای $a_i<0$ خواهیم داشت
$a_i f(\mathbf{x}_i)>0$.

پس الزاماً

\[
\sum_{i=1}^{D+2}
a_i f(\mathbf{x}_i)>0.
\]

اما از روابط وابستگی خطی نتیجه می‌شود:

\[
\begin{aligned}
\sum_{i=1}^{D+2}a_i f(\mathbf{x}_i)
&=
\mathbf{w}^{T}
\sum_{i=1}^{D+2}a_i\mathbf{x}_i
+
b\sum_{i=1}^{D+2}a_i\\
&=0,
\end{aligned}
\]

که تناقض است.

پس هیچ مجموعه $D+2$ نقطه‌ای
قابل شکستن نیست و نتیجه می‌گیریم:

\[
\boxed{
\operatorname{VCdim}(\mathcal{H}_{\mathrm{lin}})
=D+1.
}
\]

این نتیجه نشان می‌دهد که افزایش
بعد فضای ورودی، ظرفیت دسته‌بندهای
خطی آفین را افزایش می‌دهد.


% ==========================================
% 4. Growth function and Sauer-Shelah
% ==========================================

\subsection{تابع رشد و لم ساوئر و شلاه}

بعد وی‌سی تنها مشخص می‌کند
بزرگ‌ترین مجموعه قابل شکستن چه اندازه‌ای دارد.
برای شمارش تعداد برچسب‌گذاری‌های قابل تحقق
روی مجموعه‌های بزرگ‌تر، از
\textbf{تابع رشد}
\en{(Growth Function)}
استفاده می‌کنیم:

\[
\boxed{
\Pi_{\mathcal{H}}(n)
=
\max_{\substack{
S\subseteq\mathcal{X}\\
|S|=n
}}
|\mathcal{H}(S)|.
}
\]

واضح است که

\[
\Pi_{\mathcal{H}}(n)\leq2^n.
\]

اگر بعد وی‌سی کلاس برابر
$d$ باشد، برای $n\leq d$ داریم:

\[
\Pi_{\mathcal{H}}(n)=2^n.
\]

اما برای $n>d$،
همه برچسب‌گذاری‌ها قابل تحقق نیستند.

لم ساوئر و شلاه
\en{(Sauer--Shelah Lemma)}
بیان می‌کند که اگر
$\operatorname{VCdim}(\mathcal{H})=d<\infty$،
آنگاه

\[
\boxed{
\Pi_{\mathcal{H}}(n)
\leq
\sum_{i=0}^{d}\binom{n}{i}.
}
\]

به‌ویژه برای $n\geq d\geq1$ داریم:

\[
\boxed{
\Pi_{\mathcal{H}}(n)
\leq
\left(\frac{en}{d}\right)^d.
}
\]

در نتیجه، برای یک کلاس با بعد وی‌سی
ثابت و متناهی، تعداد الگوهای برچسب‌گذاری
قابل تحقق، حداکثر رشد چندجمله‌ای
نسبت به تعداد نقاط دارد؛
در حالی که تعداد تمام برچسب‌گذاری‌های
ممکن برابر $2^n$ و نمایی است.

بنابراین بعد وی‌سی، حتی برای
کلاس‌های فرضیه نامتناهی، محدودیتی
بر پیچیدگی ترکیبیاتی مدل ایجاد می‌کند.


% ==========================================
% 5. VC dimension and generalization
% ==========================================

\subsection{ارتباط بعد وی‌سی با تعمیم‌پذیری}

اهمیت آماری بعد وی‌سی از ارتباط آن
با تفاوت میان خطای آموزشی و خطای واقعی
ناشی می‌شود.

فرض کنید $N$ نمونه مستقل و هم‌توزیع
از یک توزیع ناشناخته $\mathcal{D}$
در اختیار داریم.

ریسک واقعی یک فرضیه $h$ را
در دسته‌بندی دودویی با زیان صفر و یک،
به صورت زیر تعریف می‌کنیم:

\[
R(h)=
\Pr_{(\mathbf{x},y)\sim\mathcal{D}}
\bigl(h(\mathbf{x})\neq y\bigr).
\]

در مقابل، ریسک تجربی روی مجموعه آموزشی
برابر است با

\[
\widehat{R}_N(h)=
\frac{1}{N}
\sum_{i=1}^{N}
\mathbb{I}
\bigl\{
h(\mathbf{x}_i)\neq y_i
\bigr\}.
\]

هدف، آن است که خطای آموزشی
تخمین مناسبی از خطای واقعی باشد،
حتی زمانی که مدل پس از مشاهده داده‌ها
انتخاب شده است.

یکی از نامساوی‌های استاندارد
در نظریه وی‌سی بیان می‌کند که
برای هر $\varepsilon>0$،

\[
\Pr\left(
\sup_{h\in\mathcal{H}}
|R(h)-\widehat{R}_N(h)|
>\varepsilon
\right)
\leq
8\Pi_{\mathcal{H}}(2N)
\exp\left(
-\frac{N\varepsilon^2}{32}
\right).
\]

با فرض
$d=\operatorname{VCdim}(\mathcal{H})$
و $N\geq d\geq1$،
می‌توان از لم ساوئر و شلاه استفاده کرد:

\[
\Pi_{\mathcal{H}}(2N)
\leq
\left(\frac{2eN}{d}\right)^d.
\]

در نتیجه، با احتمال دست‌کم
$1-\delta$، هم‌زمان برای تمام
$h\in\mathcal{H}$ داریم:

\[
\boxed{
|R(h)-\widehat{R}_N(h)|
\leq
\sqrt{
\frac{32}{N}
\left[
d\ln\left(\frac{2eN}{d}\right)
+
\ln\left(\frac{8}{\delta}\right)
\right]
}.
}
\]

این رابطه نشان می‌دهد که با
افزایش تعداد نمونه‌ها و ثابت ماندن
بعد وی‌سی، کران اختلاف خطای آموزشی
و واقعی کاهش می‌یابد.

از سوی دیگر، افزایش $d$،
برای ثابت نگه داشتن تضمین تعمیم،
معمولاً نیاز به تعداد بیشتری نمونه دارد.

البته این نامساوی یک کران آماری
از نوع بدترین حالت و مستقل از
توزیع داده است؛ بنابراین
تعداد دقیق نمونه‌های لازم
در هر مسئله واقعی را تعیین نمی‌کند.


% ==========================================
% 6. Connection to Question 2
% ==========================================

\subsection{ارتباط با چندجمله‌ای‌ها و نفرین ابعاد}

در پرسش دوم نشان دادیم که تعداد
ضرایب مستقل یک چندجمله‌ای
با درجه حداکثر $M$ در فضای
$\mathbb{R}^{D}$ برابر است با

\[
P(D,M)=\binom{D+M}{M}.
\]

اکنون کلاس تمام دسته‌بندهای حاصل
از تعیین علامت چنین چندجمله‌ای‌هایی
را تعریف می‌کنیم:

\[
\mathcal{H}_{D,M}
=
\left\{
h_p(\mathbf{x})
=
\mathbb{I}\{p(\mathbf{x})\geq0\}
:
\deg(p)\leq M
\right\}.
\]

در این تعریف، تمام ضرایب حقیقی
چندجمله‌ای آزاد هستند و فضای ورودی
کل $\mathbb{R}^{D}$ است.

نشان می‌دهیم که

\[
\boxed{
\operatorname{VCdim}(\mathcal{H}_{D,M})
=
\binom{D+M}{M}.
}
\]

\par\medskip\noindent\textbf{اثبات.}
فرض کنید

\[
P=\binom{D+M}{M}
\]

و $\boldsymbol{\phi}(\mathbf{x})\in\mathbb{R}^{P}$
بردار تمام تک‌جمله‌ای‌های
با درجه حداکثر $M$ باشد.
یکی از مؤلفه‌های این بردار،
تک‌جمله‌ای ثابت $1$ است.

هر چندجمله‌ای مورد نظر را می‌توان نوشت:

\[
p(\mathbf{x})
=
\boldsymbol{\theta}^{T}
\boldsymbol{\phi}(\mathbf{x}),
\qquad
\boldsymbol{\theta}\in\mathbb{R}^{P}.
\]

\subsubsection{کران بالا}

اگر مؤلفه ثابت را جدا کنیم،
این دسته‌بند به یک دسته‌بند
خطی آفین در فضای ویژگی
$\mathbb{R}^{P-1}$ تبدیل می‌شود.

از آنجا که بعد وی‌سی
دسته‌بندهای آفین در فضای
$\mathbb{R}^{P-1}$ برابر $P$ است،
محدود کردن ورودی آن‌ها به
بردارهای ویژگی حاصل از داده‌ها
نمی‌تواند بعد وی‌سی را افزایش دهد.
پس

\[
\operatorname{VCdim}(\mathcal{H}_{D,M})
\leq P.
\]

\subsubsection{کران پایین}

تک‌جمله‌ای‌های متمایز، به عنوان
توابع حقیقی روی $\mathbb{R}^{D}$،
مستقل خطی هستند.

بنابراین بردارهای

\[
\left\{
\boldsymbol{\phi}(\mathbf{x})
:
\mathbf{x}\in\mathbb{R}^{D}
\right\}
\]

فضای $\mathbb{R}^{P}$ را پوشش خطی می‌دهند.

برای اثبات این ادعا، فرض کنید
خلاف آن برقرار باشد.
در این صورت یک بردار ناصفر
$\mathbf{a}\in\mathbb{R}^{P}$
وجود خواهد داشت که

\[
\mathbf{a}^{T}
\boldsymbol{\phi}(\mathbf{x})=0
\qquad
\forall\mathbf{x}\in\mathbb{R}^{D}.
\]

اما طرف چپ یک چندجمله‌ای با
ضرایب غیرهمگی صفر است که روی
تمام $\mathbb{R}^{D}$ صفر می‌شود؛
این موضوع با استقلال خطی
تک‌جمله‌ای‌ها تناقض دارد.

پس می‌توان $P$ نقطه
$\mathbf{x}_1,\ldots,\mathbf{x}_P$
انتخاب کرد که بردارهای ویژگی
آن‌ها مستقل خطی باشند.

ماتریس ویژگی این نقاط را
به صورت زیر تعریف می‌کنیم:

\[
\boldsymbol{\Phi}=
\begin{bmatrix}
\boldsymbol{\phi}(\mathbf{x}_1)^T\\
\vdots\\
\boldsymbol{\phi}(\mathbf{x}_P)^T
\end{bmatrix}
\in\mathbb{R}^{P\times P}.
\]

این ماتریس وارون‌پذیر است.

اکنون هر برچسب‌گذاری دلخواه
$\mathbf{s}\in\{-1,+1\}^{P}$
را در نظر بگیرید و انتخاب کنید:

\[
\boldsymbol{\theta}
=
\boldsymbol{\Phi}^{-1}\mathbf{s}.
\]

در این صورت،

\[
\boldsymbol{\Phi}\boldsymbol{\theta}
=\mathbf{s},
\]

و بنابراین

\[
p(\mathbf{x}_i)=s_i,
\qquad i=1,\ldots,P.
\]

پس علامت خروجی چندجمله‌ای
با تمام برچسب‌های دلخواه منطبق است.
در نتیجه این $P$ نقطه قابل شکستن‌اند و

\[
\operatorname{VCdim}(\mathcal{H}_{D,M})
\geq P.
\]

با ترکیب دو کران، نتیجه نهایی به دست می‌آید:

\[
\boxed{
\operatorname{VCdim}(\mathcal{H}_{D,M})
=
P(D,M)
=
\binom{D+M}{M}.
}
\]

\subsubsection{ارتباط با نفرین ابعاد}
این برابری نشان می‌دهد که
در کلاس کامل دسته‌بندهای چندجمله‌ای،
افزایش تعداد پارامترها مستقیماً
به افزایش بعد وی‌سی منجر می‌شود.

برای درجه ثابت $M$،
طبق نتیجه پرسش دوم داریم:

\[
\operatorname{VCdim}(\mathcal{H}_{D,M})
\sim\frac{D^M}{M!}
\qquad(D\to\infty).
\]

بنابراین در این حالت رشد ظرفیت
چندجمله‌ای است.

اما اگر درجه همراه با بعد افزایش یابد،
برای مثال $M=D$، خواهیم داشت:

\[
\operatorname{VCdim}(\mathcal{H}_{D,D})
=
\binom{2D}{D}
\sim
\frac{4^D}{\sqrt{\pi D}},
\]

که دارای رشد نمایی نسبت به $D$ است.

این نتیجه، همراه با کران‌های
تعمیم مبتنی بر بعد وی‌سی،
نشان می‌دهد چرا افزایش ابعاد
می‌تواند یادگیری قابل اعتماد
را از نظر آماری دشوارتر کند.

البته برابری بعد وی‌سی و تعداد
پارامترها یک قاعده عمومی برای
تمام مدل‌های یادگیری ماشین نیست؛
بلکه در اینجا برای کلاس مشخص
دسته‌بندهای چندجمله‌ای آزاد
اثبات شده است.


% ==========================================
% 7. Conclusion
% ==========================================

\subsection{جمع‌بندی}

بعد وی‌سی، بیشترین تعداد نقاطی است
که یک کلاس فرضیه می‌تواند تمام
برچسب‌گذاری‌های دودویی ممکن روی
آن‌ها را محقق کند.

این مفهوم از یک سو ظرفیت ترکیبیاتی
مدل را اندازه‌گیری می‌کند و از سوی دیگر،
از طریق نظریه تعمیم آماری،
رابطه ظرفیت مدل با تعداد نمونه‌های
آموزشی را روشن می‌سازد.

بنابراین اهمیت بعد وی‌سی صرفاً
در توانایی برازش داده‌های آموزشی نیست،
بلکه در فراهم کردن چارچوبی ریاضی
برای بررسی تعادل میان پیچیدگی مدل،
میزان داده و تعمیم‌پذیری آن است.

\end{document}
